\subsubsection{Matériel utilisé}

\subsubsubsection{Matériaux}
\begin{longtable}{|p{2cm}|p{4cm}|p{3cm}|p{3cm}|p{3cm}|}
    \hline
    \textbf{Type} & \textbf{Utilisation} & \textbf{Avantages} & \textbf{Inconvénients} & \textbf{Recommandations} \\
    \hline
    Fibre de verre & Tubes, ailerons, coiffe & Moins cher, facile à usiner & Plus lourd que le carbone & À utiliser pour les pièces non critiques. \\
    \hline
    Fibre de carbone & Renforts locaux, protection thermique & Légèreté, résistance & Coût élevé, fabrication complexe & Privilégier les tissus 200 g/m² pour un bon compromis résistance/poids. \\
    \hline
    Aluminium 7075 \& 6061 & Fixations, embases moteur & Résistant, usinable & Lourd & Limiter à les pièces soumises à des efforts élevés. \\
    \hline
    Colles époxy & Assemblage des composites & Bonne adhérence & Sensible à la température & Utiliser des colles haute température (ex. : Araldite 2011). \\
    \hline
    \caption{Matériaux utilisés pour la fusée SP-02}
    \label{tab:materiaux}
\end{longtable}

\subsubsubsection{Outillage}
\begin{longtable}{|p{4cm}|p{4.5cm}|p{6.5cm}|}
    \hline
    \textbf{Outil} & \textbf{Utilisation} & \textbf{Retours} \\
    \hline
    Gabarits PLA & Maintien des ailerons pour collage, fabrication de pièces composites tels cages parachutes & Précision moyenne (impression 3D) \\
    \hline
    Imprimante 3D & Prototypes, fixations électroniques & Rapide pour les pièces complexes, mais résistance limitée pour les pièces structurelles. \\
    \hline
    Machine CNC & Usinage des pièces structurelles en aluminium & Précision élevée nécessaire pour les tolérances serrées. \\
    \hline
    \caption{Outillage utilisé pour la fabrication de la fusée SP-02}
    \label{tab:outillage}
\end{longtable}

\subsubsubsection{Logiciels}
\begin{longtable}{|p{2cm}|p{5cm}|p{6cm}|}
    \hline
    \textbf{Logiciel} & \textbf{Utilisation} & \textbf{Retours} \\
    \hline
    OpenRocket & Simulation de vol & Simple d’utilisation, mais limité pour les fusées bi-étages. \\
    \hline
    Patran & Simulations mécaniques (FEM) & Puissance limitée par la licence étudiante \\
    \hline
    StarCCM+ & Simulations d'écoulement de fluide (CFD) & Puissance limitée par la licence étudiante \\
    \hline
    EasyEDA & Conception des PCB & Open-source et efficace, optimisé avec l'usineur JLC \\
    \hline
    CATIA & Modélisation 3D & Idéal pour l’intégration mécanique/électronique. \\
    \hline
    CubeMX & Programmation des microconôcoleurs/debuggage & Pratique \\
    \hline
    VSCode & Codage des programmes et rapport de projet & Optimisé \\
    \hline
    Stabtraj & Modélisation de la fusée sur Excel & Très Pratique \\
    \hline
    \caption{Logiciels utilisés pour le projet SP-02}
    \label{tab:logiciels}
\end{longtable}

\subsubsection*{Recommandations pour les futurs projets}
\begin{itemize}
    \item \textbf{Matériaux} : Tester des composites hybrides (carbone/kevlar) et des formes de moules variées pour accroître les possiblités de conception et de réalisation.
    \item \textbf{Outillage} : Investir dans des pinces à sertir de qualité.
    \item \textbf{Électronique} : Utiliser des modules certifiés spatial (ex. : capteurs TE Connectivity) pour une meilleure fiabilité.
    \item \textbf{Logiciels} : Former les équipes en amont du début du projet pour gagner du temps.
\end{itemize}
